\subsection{Justificación pedagógica de las relaciones}
\vspace{4pt}

\begin{tablaSAlarga}{colspec={Q[l,wd=1cm] Q[l,wd=3.4cm] Q[c,wd=1.4cm] Q[l,wd=8.9cm] Q[c,wd=1.3cm] Q[l,wd=7cm]},
  rowsep=2pt, colsep=4pt, cells={font=\scriptsize}, row{1}={bg=urjc, fg=white, font=\footnotesize\bfseries},}
Crit. & Aspecto & Sesiones & Por qué en esas sesiones & Instrum. & Por qué ese instrumento \\
\textbf{7.1} & Representar y operar con vectores & 2, 3, 4 & Se construye primero en F1 (S2), para que cada operación tenga un significado geométrico antes de formalizarla en S3 y S4. & CU, FI & La ficha recoge la representación con tecnología y Q1 comprueba su interpretación. El examen escrito no permite usar GeoGebra. \\
\textbf{2.1} & Validar cálculos vectoriales & 2, 4, 18 & Validar exige tener un cálculo propio: en S2 se contrasta de forma exploratoria y en S4, con la fórmula de la proyección, se explica la discrepancia. S18 comprueba si se hace sin herramienta. & EX, FI & F1 registra el contraste entre mano y GeoGebra; la pregunta 2 del examen obliga a detectar un error ajeno. \\
\textbf{5.2} & Producto escalar en problemas & 4, 18 & El producto escalar une coordenadas y geometría, y se enseña en S4; S18 comprueba que se aplica de forma autónoma. & EX, CU & Q1 da un diagnóstico temprano para el refuerzo por criterio de S13, y el examen verifica la transferencia. \\
\textbf{3.2} & Investigar con la circunferencia & 5 & S5 precede a propósito a la exposición de S6: hay que investigar antes de conocer la fórmula. & FI, OB & La ficha muestra lo investigado; la observación capta si se explora antes de pedir ayuda, que no queda en el papel. \\
\textbf{3.1} & Conjeturar y demostrar & 5, 6, 18 & La conjetura nace en S5 y se demuestra en S6, de modo que la demostración responde a una pregunta propia. S18 comprueba que se demuestra sin apoyo. & EX, FI & F2 recoge la conjetura y la pregunta 3a del examen, la demostración individual. \\
\textbf{1.2} & Ecuaciones trigonométricas & 7, 18 & Las ecuaciones usan el ángulo doble, que se estudia en la misma sesión, y la pieza del reto les da un uso inmediato. & EX, CU & Obtener todas las soluciones es un procedimiento: Q2 da retorno inmediato y el examen lo certifica. \\
\textbf{1.1} & Elegir el teorema adecuado & 8, 9, 18 & La elección se aprende en S8 con datos variados y se pone a prueba en S9, con el caso que genera más dudas. & EX, FI & F3 y la pregunta 4 piden justificar la elección, no solo dar el resultado. \\
\textbf{1.2} & Caso ambiguo & 9, 18 & Va al final del bloque porque exige dominar antes el teorema del seno. & EX, FI & En F3 se ve cuántas soluciones hay al arrastrar; la pregunta 4b exige obtenerlas por cálculo. \\
\textbf{2.1} & Validar triángulos & 9 & S9 concentra el caso en que validar es decisivo: si $\sin B>1$, el triángulo no existe. & CU, FI & La comprobación arrastrando vértices solo se ve en F3; Q3 la cierra con un caso sin solución. \\
\textbf{5.1} & Conectar vector y recta & 10, 13, 18 & S10 deduce las ecuaciones del vector director; el mapa llega en S13, cuando ya se han visto los tres bloques y hay algo que conectar. & EX, MC & El mapa hace visibles las conexiones y la pregunta 5 del examen comprueba que se usan. \\
\textbf{5.2} & Ángulos y distancias & 11, 18 & Reutilizan el producto escalar de S4; por eso se sitúan al final del bloque, para volver sobre él. & EX, CU & Q4 alimenta el refuerzo por criterio de S13 y el examen verifica el dominio individual. \\
\textbf{7.2} & Elegir la ecuación de la recta & 10, 12, 18 & S10 presenta las cinco formas y F4, en S12, obliga a elegir una para un problema real. & EX, FI & F4 y la pregunta 5d piden explicar por qué se elige esa ecuación. \\
\textbf{1.1} & Método y modelo & 14, 15, 19, 20 & Se aplica cuando ya se dispone de todas las herramientas: el método se aprueba en S13, se ejecuta en S14 y se modela en S15. & LI, EP & El libro documenta el proceso; la defensa comprueba que todo el equipo lo entiende. \\
\textbf{2.2} & Elegir el resultado & 14, 15, 19, 20 & Elegir solo tiene sentido con medidas reales y repetidas (S14), y la sensibilidad al error se analiza con ellas en S15. & LI, EP & El libro recoge datos y decisión; en la defensa hay que justificarla ante preguntas. \\
\textbf{7.2} & Integrar representaciones & 15, 16, 17, 19, 20 & Integrar exige tener ya datos y modelo (S15 y S16), y S17 permite mejorarlo; las defensas muestran si otras personas lo entienden. & LI, CO & El libro es la representación. La coevaluación aporta lo que pide el criterio: su utilidad para compartir información. \\
\textbf{8.1} & Comunicar con rigor & 17, 19, 20 & Se ensaya en S17 y se evalúa en S19 y S20, ante un público real. & EP, CO & La comunicación oral solo se puede evaluar en una exposición; la co/autoevaluación añade la mirada del público. \\
\textbf{9.1} & Aprender del error & 5, 12, 14, 19, 20 & Se observa donde el error aparece de forma natural: conjeturas falsas (S5), construcciones (S12) y medidas (S14). & OB, CO & Una actitud no cabe en un examen: se observa cuando ocurre y se reflexiona en la autoevaluación. \\
\textbf{9.2} & Aceptar la crítica & 16, 17, 19, 20 & Necesita una crítica aprovechable: la coevaluación del borrador en S16 se aprovecha en S17, dentro del horario de clase. & OB, CO & El registro de cambios permite ver qué se ha incorporado realmente. \\
\textbf{9.3} & Colaborar en equipo & 2, 14, 15, 16, 19, 20 & Los roles se presentan en S2 y se exigen sobre todo de S14 a S16, cuando el producto depende de todos. & OB, CO & Solo la observación y los compañeros pueden valorar la colaboración efectiva. \\
\textbf{6.1} & Conectar con la instalación & 14, 16, 19, 20 & La conexión con el criterio técnico requiere haber medido; se hace en S16, con la orientación del departamento de Tecnología. & LI, EP & La propuesta recoge el cálculo y la defensa, su viabilidad ante quien decide. \\
\textbf{6.2} & Valorar impacto y límites & 16, 19, 20 & La reflexión tiene sentido cuando el modelo ya existe y se conocen sus resultados. & LI, EP & El apartado del libro la deja por escrito y la defensa la somete a preguntas. \\
\end{tablaSAlarga}
